\documentclass[12pt]{article}
\usepackage[utf8]{inputenc}\usepackage[T1]{fontenc}
\usepackage[a4paper,margin=2.5cm]{geometry}
\usepackage{amsmath,amssymb,amsthm,mathtools,mathrsfs}
\usepackage[dvipsnames]{xcolor}
\usepackage{booktabs,enumitem,hyperref}
\setlist{nosep,leftmargin=1.8em}
\hypersetup{colorlinks=true,linkcolor=blue!60!black,citecolor=green!40!black,urlcolor=blue!50!black}

\theoremstyle{plain}
\newtheorem{theorem}{Theorem}[section]
\newtheorem{proposition}[theorem]{Proposition}
\newtheorem{corollary}[theorem]{Corollary}
\newtheorem{lemma}[theorem]{Lemma}
\theoremstyle{definition}
\newtheorem{definition}[theorem]{Definition}
\newtheorem{axiom}[theorem]{Axiom}
\newtheorem{example}[theorem]{Example}
\newtheorem{construction}[theorem]{Construction}
\theoremstyle{remark}
\newtheorem{remark}[theorem]{Remark}

\newcommand{\Aut}{\mathrm{Aut}}
\newcommand{\End}{\mathrm{End}}
\newcommand{\Ker}{\mathrm{Ker}}
\newcommand{\Def}{\mathrm{Def}}
\newcommand{\Stab}{\mathrm{Stab}}
\newcommand{\Isom}{\mathrm{Isom}}
\newcommand{\Twin}{\mathrm{Twin}}
\newcommand{\Poinc}{\mathrm{Poinc}}
\newcommand{\Diff}{\mathrm{Diff}}
\newcommand{\SO}{\mathrm{SO}}
\newcommand{\SU}{\mathrm{SU}}
\newcommand{\SL}{\mathrm{SL}}
\newcommand{\GL}{\mathrm{GL}}
\newcommand{\UU}{\mathrm{U}}
\newcommand{\ISO}{\mathrm{ISO}}
\newcommand{\dd}{\mathrm{d}}
\newcommand{\CC}{\mathbb{C}}
\newcommand{\RR}{\mathbb{R}}
\newcommand{\ZZ}{\mathbb{Z}}
\newcommand{\HH}{\mathbb{H}}

\title{\textbf{The Theory of Twin Spaces:\\General Framework, Real Spacetime,\\and Complex Spacetime with Field Interactions}}

\author{Jocelyn Nemb\'e\thanks{Email: jnembe777@gmail.com, jnembe@root-insights.org}\\[0.3em]
\small Roots Insights Laboratory, Singapore\\
\small National Institute of Management Sciences, Libreville, Gabon}

\date{}

\begin{document}
\maketitle

% ============================================================================
\begin{abstract}
We present a self-contained exposition of the theory of twin spaces, from its foundations on arbitrary sets to its application to real and complex spacetimes with gauge interactions. \textbf{Part~I} develops the general theory: the generation of twin operations from a base operation via gauge maps, the kernel-deformation splitting, equivariant maps between $(G,H)$-spaces, isomorphism theorems, and orbit-class theorems. \textbf{Part~II} applies the framework to the real Minkowski spacetime $\RR^{1,3}$, recovering Lorentz transformations, Einstein addition, and the Poincar\'e group as the kernel. \textbf{Part~III} constructs a \emph{complex spacetime} $\CC^{1,3}$ equipped with the complex twin hierarchy $z_1 \odot_n z_2 = \exp^{(n)}(\log^{(n)}(z_1) + \log^{(n)}(z_2))$, where $\exp^{(n)}$ and $\log^{(n)}$ denote iterated complex exponential and logarithm. The complex hierarchy generates a countably infinite family of spacetime arithmetics parameterized by $n \in \ZZ$, each with its own isometry group. We construct a \textbf{gauge-gravity model} on $\CC^{1,3}$ by crossing the complex twin group with gauge groups representing interactions ($\SU(3) \times \SU(2) \times \UU(1)$), producing a unified framework where gravity (from the twin hierarchy) and gauge forces (from the fiber) coexist in a single $(G,H)$-equivariant structure.
\end{abstract}

\tableofcontents
\newpage


% ============================================================================
% PART I: GENERAL THEORY
% ============================================================================

\part{General Theory of Twin Spaces}

% ============================================================================
\section{Twin Operations on an Arbitrary Set}
\label{sec:generation}
% ============================================================================

\subsection{The generation mechanism}

\begin{definition}[Base algebraic structure]\label{def:base}
Let $(E, \alpha)$ be a set $E$ equipped with a binary operation $\alpha: E \times E \to E$.
\end{definition}

\begin{definition}[Gauge map]\label{def:gauge}
A \textbf{gauge map} is a bijection $\phi: E \to E$. The set of all gauge maps forms the symmetric group $\mathrm{Bij}(E)$.
\end{definition}

\begin{definition}[Twin operation]\label{def:twin-op}
Given $(E, \alpha)$ and a gauge map $\phi$, the \textbf{twin operation} $\odot_\phi$ is:
\begin{equation}\label{eq:twin-def}
\boxed{x \odot_\phi y := \phi^{-1}(\phi(x) \;\alpha\; \phi(y))}
\end{equation}
Equivalently: conjugate $\alpha$ by $\phi$. The operation $\odot_\phi$ inherits all equational properties of $\alpha$ (associativity, commutativity, existence of identity/inverse) because conjugation preserves equations.
\end{definition}

\begin{theorem}[Property inheritance]\label{thm:inheritance}
If $(E, \alpha)$ is a group (resp.\ ring, field, vector space), then $(E, \odot_\phi)$ is a group (resp.\ ring, field, vector space) with:
\begin{itemize}
    \item Identity: $e_\phi = \phi^{-1}(e_\alpha)$ where $e_\alpha$ is the identity of $\alpha$.
    \item Inverse: $x^{-1}_\phi = \phi^{-1}(\phi(x)^{-1}_\alpha)$.
    \item Homomorphism: $\phi: (E, \odot_\phi) \to (E, \alpha)$ is an isomorphism.
\end{itemize}
\end{theorem}

\begin{proof}
For associativity: $x \odot_\phi (y \odot_\phi z) = \phi^{-1}(\phi(x) \;\alpha\; \phi(\phi^{-1}(\phi(y) \;\alpha\; \phi(z)))) = \phi^{-1}(\phi(x) \;\alpha\; \phi(y) \;\alpha\; \phi(z)) = (x \odot_\phi y) \odot_\phi z$. The other properties follow similarly. The map $\phi: (E, \odot_\phi) \to (E, \alpha)$ satisfies $\phi(x \odot_\phi y) = \phi(x) \;\alpha\; \phi(y)$ by definition, so it is a homomorphism; being a bijection, it is an isomorphism.
\end{proof}

\subsection{The twin hierarchy}

\begin{definition}[Twin hierarchy]\label{def:hierarchy}
Given $(E, \alpha)$ and a fixed gauge map $\phi$, the \textbf{twin hierarchy} is the family of operations $\{\odot_n\}_{n \in \ZZ}$ defined by:
\begin{equation}
x \odot_n y := (\phi^{-1})^{(n)}(\phi^{(n)}(x) \;\alpha\; \phi^{(n)}(y))
\end{equation}
where $\phi^{(n)}$ denotes the $n$-fold composition of $\phi$ (with $\phi^{(0)} = \mathrm{id}$, $\phi^{(-n)} = (\phi^{-1})^{(n)}$). Note $\odot_0 = \alpha$ and $\odot_1 = \odot_\phi$.
\end{definition}

\begin{example}[The arithmetic hierarchy on $\RR_{>0}$]\label{ex:arithmetic}
Take $E = \RR_{>0}$, $\alpha = +$ (addition), gauge map $\phi = \log$ (so $\phi^{-1} = \exp$ and $\odot_n(x,y) = \exp^{(n)}(\log^{(n)} x + \log^{(n)} y)$, consistent with the complex hierarchy of \S\ref{sec:complex-hierarchy}):
\begin{center}\small
\begin{tabular}{cll}\toprule
$n$ & $\odot_n$ & \textbf{Name} \\\midrule
$-1$ & $x \odot_{-1} y = \ln(e^x + e^y)$ & LogSumExp \\
$0$ & $x \odot_0 y = x + y$ & Addition \\
$1$ & $x \odot_1 y = x \cdot y$ & Multiplication \\
$2$ & $x \odot_2 y = x^{\ln y}$ & Power tower \\
\bottomrule
\end{tabular}
\end{center}
Each level generates the next via $\exp$-conjugation. The hierarchy extends in both directions.
\end{example}

\begin{example}[The Lorentz hierarchy on $(-c, c)$]\label{ex:lorentz}
Take $E = (-c, c)$, $\alpha = +$, $\phi = \mathrm{artanh}(\cdot/c)$:
$v_1 \odot_1 v_2 = c\tanh(\mathrm{artanh}(v_1/c) + \mathrm{artanh}(v_2/c)) = \frac{v_1 + v_2}{1 + v_1v_2/c^2}$.
This is Einstein velocity addition.
\end{example}


% ============================================================================
\section{The Kernel-Deformation Splitting}
\label{sec:kernel}
% ============================================================================

\begin{definition}[Group action on operations]\label{def:group-action}
Let $G$ be a group acting on $(E, \alpha)$ by bijections: for each $g \in G$, $g: E \to E$ is a bijection. Each $g$ generates a twin operation $\odot_g$ via Definition~\ref{def:twin-op}. The map $g \mapsto \odot_g$ partitions $G$ into classes.
\end{definition}

\begin{definition}[Kernel]\label{def:kernel}
The \textbf{kernel} of the action is the stabilizer of the base operation:
\begin{equation}
\Ker(\alpha) := \{g \in G : \odot_g = \alpha\} = \{g \in G : g(x \;\alpha\; y) = g(x) \;\alpha\; g(y) \;\forall x, y\}
\end{equation}
Elements of $\Ker(\alpha)$ are the \textbf{automorphisms} of $(E, \alpha)$: they preserve the operation. In physics, these are the \textbf{symmetries}.
\end{definition}

\begin{definition}[Deformation space]\label{def:deformation}
The \textbf{deformation space} (or twin space) is the quotient:
\begin{equation}
\Twin(\alpha) := G / \Ker(\alpha)
\end{equation}
It parameterizes the \emph{physically distinct} twin operations: $g_1$ and $g_2$ produce the same operation iff $g_1 g_2^{-1} \in \Ker(\alpha)$.
\end{definition}

\begin{theorem}[The splitting theorem]\label{thm:splitting}
The Lie algebra $\mathfrak{g}$ of $G$ admits a splitting:
\begin{equation}\label{eq:splitting}
\boxed{\mathfrak{g} = \underbrace{\mathfrak{ker}(\alpha)}_{\text{symmetries}} \oplus \underbrace{\mathfrak{def}(\alpha)}_{\text{deformations}}}
\end{equation}
where $\mathfrak{ker}(\alpha) = \{\xi \in \mathfrak{g} : \mathcal{L}_\xi\alpha = 0\}$ (infinitesimal automorphisms) and $\mathfrak{def}(\alpha)$ is a complement. The kernel dimension $\kappa := \dim(\mathfrak{ker})$ counts the independent symmetries.
\end{theorem}

\begin{proof}
$\mathfrak{ker}(\alpha) = \ker(\xi \mapsto \mathcal{L}_\xi\alpha)$ is the kernel of a linear map from $\mathfrak{g}$ to the space of bilinear maps on $E$. It is a Lie subalgebra (since $[\xi_1, \xi_2]$ preserves $\alpha$ if both $\xi_i$ do). Any complement $\mathfrak{def}$ completes the splitting.
\end{proof}

\begin{theorem}[Kernel characterization]\label{thm:kernel-char}
The kernel $\Ker(\alpha)$ is characterized by:
\begin{enumerate}[label=(\alph*)]
    \item $\Ker(\alpha)$ is a closed subgroup of $G$ (intersection of automorphism conditions).
    \item $\Ker(\alpha) = \bigcap_{x,y \in E} \mathrm{Stab}(\alpha(x,y))$ (stabilizer of all products).
    \item If $E$ is finite-dimensional, $\kappa \leq \dim(G)$ with equality iff $G$ acts by automorphisms.
    \item For a metric $g$ on a manifold: $\Ker(g) = \Isom(\mathcal{M}, g)$ (isometry group), and its Lie algebra $\mathfrak{ker}(g)$ consists exactly of the Killing vector fields.
\end{enumerate}
\end{theorem}


% ============================================================================
\section{Equivariant Maps Between Twin Spaces}
\label{sec:equivariant}
% ============================================================================

\begin{definition}[$(G,H)$-equivariant map]\label{def:equivariant}
Let $(E_1, \alpha_1, G)$ and $(E_2, \alpha_2, H)$ be two spaces with group actions. A \textbf{$(G,H)$-equivariant map} is a pair $(\Phi, \varphi)$ where $\Phi: E_1 \to E_2$ is a map and $\varphi: G \to H$ is a group homomorphism, satisfying:
\begin{equation}\label{eq:equivariant}
\boxed{\Phi(g \cdot x) = \varphi(g) \cdot \Phi(x) \qquad \forall g \in G,\; x \in E_1}
\end{equation}
The map $\Phi$ ``intertwines'' the actions of $G$ and $H$ via $\varphi$.
\end{definition}

\begin{theorem}[Kernel of equivariant maps]\label{thm:kernel-equivariant}
If $(\Phi, \varphi)$ is $(G,H)$-equivariant, then:
\begin{enumerate}[label=(\alph*)]
    \item $\ker(\varphi) \subseteq \Ker_\Phi := \{g \in G : \Phi(g \cdot x) = \Phi(x)\;\forall x\}$.
    \item $\Phi$ maps $G$-orbits to $H$-orbits: $\Phi(\mathcal{O}_G(x)) \subseteq \mathcal{O}_H(\Phi(x))$.
    \item $\Phi$ maps $G$-invariants to $H$-invariants.
    \item If $\Phi$ is bijective and $\varphi$ is an isomorphism: $(E_1, \odot_g) \cong (E_2, \odot_{\varphi(g)})$ for all $g$.
\end{enumerate}
\end{theorem}

\begin{proof}
(a) If $g \in \ker(\varphi)$: $\Phi(g \cdot x) = \varphi(g) \cdot \Phi(x) = e_H \cdot \Phi(x) = \Phi(x)$.
(b) $\Phi(g \cdot x) = \varphi(g) \cdot \Phi(x) \in \mathcal{O}_H(\Phi(x))$.
(c) If $f(g \cdot x) = f(x)$ for all $g$, then $f \circ \Phi$ is $H$-invariant by the intertwining.
(d) Direct from the definitions.
\end{proof}

\begin{definition}[Twin morphism]\label{def:twin-morphism}
A \textbf{twin morphism} from $(E_1, \{\odot^1_g\}_{g \in G})$ to $(E_2, \{\odot^2_h\}_{h \in H})$ is a $(G,H)$-equivariant map $(\Phi, \varphi)$ such that for every $g \in G$:
\begin{equation}
\Phi(x \odot^1_g y) = \Phi(x) \odot^2_{\varphi(g)} \Phi(y)
\end{equation}
The map $\Phi$ preserves the twin operations, with the level shifted by $\varphi$.
\end{definition}


% ============================================================================
\section{Isomorphism Theorems}
\label{sec:isomorphisms}
% ============================================================================

\begin{theorem}[First isomorphism theorem for twin spaces]\label{thm:iso1}
Let $(\Phi, \varphi): (E_1, G) \to (E_2, H)$ be a surjective twin morphism. Then:
\begin{equation}
E_1 / \ker(\Phi) \cong E_2 \qquad \text{as twin spaces}
\end{equation}
where $\ker(\Phi) = \{x \in E_1 : \Phi(x) = e_2\}$ and the twin operations on the quotient are:
$[x] \odot_{[g]} [y] := [\Phi^{-1}(\Phi(x) \odot^2_{\varphi(g)} \Phi(y))]$.
\end{theorem}

\begin{proof}
Standard: $\bar\Phi: E_1/\ker(\Phi) \to E_2$ defined by $\bar\Phi([x]) = \Phi(x)$ is well-defined (since $\Phi$ is constant on cosets of $\ker(\Phi)$), bijective (surjectivity of $\Phi$ + definition of kernel), and a twin morphism (inheriting the equivariance from $\Phi$).
\end{proof}

\begin{theorem}[Second isomorphism theorem]\label{thm:iso2}
Let $N \trianglelefteq G$ be a normal subgroup contained in $\Ker(\alpha)$. Then the twin hierarchy on $E$ descends to the quotient group $G/N$:
\begin{equation}
(E, \{\odot_g\}_{g \in G}) \cong (E, \{\odot_{[g]}\}_{[g] \in G/N})
\end{equation}
In particular, twin operations indexed by elements in the same $N$-coset are identical.
\end{theorem}

\begin{theorem}[Third isomorphism theorem --- lattice correspondence]\label{thm:iso3}
There is a bijection between:
\begin{enumerate}[label=(\roman*)]
    \item Subgroups of $G$ containing $\Ker(\alpha)$, and
    \item Subgroups of $\Twin(\alpha) = G/\Ker(\alpha)$.
\end{enumerate}
This bijection preserves normality, indices, and quotients: if $\Ker(\alpha) \subseteq K \subseteq G$, then $K/\Ker(\alpha)$ is a subgroup of $\Twin(\alpha)$ and $G/K \cong \Twin(\alpha)/(K/\Ker(\alpha))$.
\end{theorem}


% ============================================================================
\section{Orbit-Class Theorems}
\label{sec:classes}
% ============================================================================

\begin{theorem}[Orbit-stabilizer for twin operations]\label{thm:orbit-stab}
For $g \in G$ acting on $(E, \alpha)$: the orbit of $\alpha$ under the subgroup generated by $g$ has size:
$|\mathcal{O}_{\langle g \rangle}(\alpha)| = [\langle g \rangle : \Stab_{\langle g \rangle}(\alpha)]$.
More generally, the number of distinct twin operations in $\{\odot_g : g \in G\}$ is $[G : \Ker(\alpha)]$.
\end{theorem}

\begin{theorem}[Burnside class counting]\label{thm:burnside}
If $G$ is finite, the number of distinct twin operations is:
\begin{equation}
|\{\odot_g : g \in G\}| = [G : \Ker(\alpha)] = \frac{|G|}{|\Ker(\alpha)|}
\end{equation}
\end{theorem}

\begin{theorem}[Classification of twin structures]\label{thm:classification}
As abstract algebraic structures all twin operations are isomorphic to $(E,\alpha)$ (Theorem~\ref{thm:inheritance}), so abstract isomorphism does not distinguish them. The meaningful equivalence is relabeling by the symmetry group: $\odot_{g_1}$ and $\odot_{g_2}$ are equivalent up to a kernel symmetry iff $g_2 \in \Ker(\alpha)\, g_1\, \Ker(\alpha)$, i.e.\ iff $g_1$ and $g_2$ lie in the same double coset $\Ker(\alpha) \backslash G / \Ker(\alpha)$. The number of twin operations inequivalent under the $\Ker(\alpha)$-action is the number of double cosets.
\end{theorem}

\begin{remark}
The plain count of \emph{distinct} operations is $[G:\Ker(\alpha)]$ (Theorem~\ref{thm:orbit-stab}); the double-coset count is the coarser enumeration after quotienting by the residual conjugation action of the symmetry group $\Ker(\alpha)$ on the set of operations.
\end{remark}


% ============================================================================
% PART II: APPLICATION TO REAL SPACETIME
% ============================================================================

\part{Application to Real Spacetime $\RR^{1,3}$}

% ============================================================================
\section{Minkowski Space as Twin Space}
\label{sec:minkowski}
% ============================================================================

\begin{construction}[Real spacetime]\label{constr:real}
Take:
\begin{itemize}
    \item $E = \RR^{1,3}$ (4-dimensional real vector space).
    \item $\alpha = \eta$ (Minkowski metric: $\eta = \mathrm{diag}(-c^2, 1, 1, 1)$, used as a bilinear form).
    \item $G = \Diff(\RR^{1,3})$ (diffeomorphism group, or $\GL(4,\RR)$ for linear transformations).
\end{itemize}
\end{construction}

\begin{theorem}[Kernel = Poincar\'e]\label{thm:poincare-kernel}
The kernel of the Minkowski metric is the Poincar\'e group:
\begin{equation}
\Ker(\eta_c) = \Poinc(c) = \SO^+(1,3) \ltimes \RR^4
\end{equation}
with $\kappa = \dim(\Poinc) = 10$. This is maximal: $\kappa = n(n+1)/2$ for $n = 4$.
\end{theorem}

\begin{theorem}[Twin operations on velocities]\label{thm:velocity-twin}
On the velocity space $(-c, c)$, the gauge map $\phi = \mathrm{artanh}(v/c)$ generates the twin operation:
\begin{equation}
v_1 \oplus_\phi v_2 = \frac{v_1 + v_2}{1 + v_1v_2/c^2}
\end{equation}
This is Einstein velocity addition. The kernel of $\oplus_\phi$ is the identity (trivial for the abelian 1D case). The deformation space $\Twin(\eta_c) = \Diff(\RR^{1,3})/\Poinc(c)$ parameterizes all possible metrics (= gravitational fields) on $\RR^{1,3}$.
\end{theorem}

\begin{remark}[GR as the deformation space]
General relativity is the physics of the deformation space $\Twin(\eta_c)$: a metric $g \neq \eta$ corresponds to a non-trivial twin operation $\odot_g \neq \odot_\eta$, and the Einstein equations constrain which elements of $\Twin(\eta_c)$ are physically realized.
\end{remark}


% ============================================================================
% PART III: COMPLEX SPACETIME
% ============================================================================

\part{Complex Spacetime $\CC^{1,3}$ and the Complex Twin Hierarchy}

% ============================================================================
\section{The Complex Twin Hierarchy}
\label{sec:complex-hierarchy}
% ============================================================================

\subsection{Complex logarithm and exponential}

\begin{definition}[Complex iterated logarithm and exponential]\label{def:complex-iter}
For $z \in \CC$ with $z \neq 0$, define:
\begin{align}
\log^{(0)}(z) &:= z, \qquad \log^{(n)}(z) := \log(\log^{(n-1)}(z)) \quad (n \geq 1), \\
\exp^{(0)}(z) &:= z, \qquad \exp^{(n)}(z) := \exp(\exp^{(n-1)}(z)) \quad (n \geq 1),
\end{align}
where $\log$ is the principal branch of the complex logarithm ($\mathrm{Im}(\log z) \in (-\pi, \pi]$).

For negative $n$: $\log^{(-n)} := \exp^{(n)}$ and $\exp^{(-n)} := \log^{(n)}$.
\end{definition}

\begin{definition}[Complex twin hierarchy]\label{def:complex-twin}
The \textbf{complex twin hierarchy} on a suitable domain $\mathcal{D}_n \subseteq \CC$ is:
\begin{equation}\label{eq:complex-twin}
\boxed{z_1 \odot_n z_2 := \exp^{(n)}(\log^{(n)}(z_1) + \log^{(n)}(z_2))}
\end{equation}
where the domain $\mathcal{D}_n$ is the maximal open set on which $\log^{(n)}$ is defined (depends on $n$ and the branch choices).
\end{definition}

\begin{theorem}[Properties of the complex hierarchy]\label{thm:complex-props}
For each $n \in \ZZ$:
\begin{enumerate}[label=(\alph*)]
    \item $(\mathcal{D}_n, \odot_n)$ is a commutative group with identity $e_n = \exp^{(n)}(0)$ and inverse $z^{-1}_n = \exp^{(n)}(-\log^{(n)}(z))$.
    \item The gauge map $\phi_n := \log^{(n)}: (\mathcal{D}_n, \odot_n) \to (\CC, +)$ is an isomorphism.
    \item Consecutive levels are related by: $z_1 \odot_{n+1} z_2 = \exp(w_1 \odot_n w_2)$ where $w_i = \log(z_i)$.
    \item The hierarchy extends the real arithmetic hierarchy: on $\RR_{>0}$, $\odot_0 = +$, $\odot_1 = \times$, and the higher levels give hyperoperations.
\end{enumerate}
\end{theorem}

\begin{proof}
(a) $\odot_n$ is conjugate to $+$ via $\log^{(n)}$, so it inherits all group properties. (b) By definition, $\log^{(n)}(z_1 \odot_n z_2) = \log^{(n)}(z_1) + \log^{(n)}(z_2)$. (c) $z_1 \odot_{n+1} z_2 = \exp^{(n+1)}(\log^{(n+1)}(z_1) + \log^{(n+1)}(z_2)) = \exp(\exp^{(n)}(\log^{(n)}(\log(z_1)) + \log^{(n)}(\log(z_2)))) = \exp(\log(z_1) \odot_n \log(z_2))$. (d) Restriction to $\RR_{>0}$ with $\log = \ln$: $\odot_0 = +$, $\odot_1 = \cdot$, as in Example~\ref{ex:arithmetic}.
\end{proof}

\begin{example}[First levels of the complex hierarchy]\label{ex:complex-levels}
\begin{center}\small
\begin{tabular}{cll}\toprule
$n$ & $z_1 \odot_n z_2$ & \textbf{Description} \\\midrule
$0$ & $z_1 + z_2$ & Complex addition \\
$1$ & $\exp(\log z_1 + \log z_2) = z_1 \cdot z_2$ & Complex multiplication \\
$2$ & $\exp(\exp(\log\log z_1 + \log\log z_2))$ & Complex hyperoperation \\
$-1$ & $\log(\exp z_1 + \exp z_2) = \mathrm{LogSumExp}(z_1, z_2)$ & Complex soft-max \\
\bottomrule
\end{tabular}
\end{center}
\end{example}


% ============================================================================
\section{Complex Spacetime $\CC^{1,3}$}
\label{sec:complex-spacetime}
% ============================================================================

\begin{definition}[Complex Minkowski spacetime]\label{def:complex-mink}
The \textbf{complex Minkowski spacetime} is:
\begin{equation}
\CC^{1,3} := \{Z = (z^0, z^1, z^2, z^3) : z^\mu \in \CC\}
\end{equation}
equipped with the complexified Minkowski metric:
\begin{equation}
\eta_\CC(Z, W) := -c^2 z^0 w^0 + z^1 w^1 + z^2 w^2 + z^3 w^3
\end{equation}
(no complex conjugation --- this is a $\CC$-bilinear form, not a Hermitian form). The real dimension is $8$; the complex dimension is $4$.
\end{definition}

\begin{remark}[Physical interpretation]
Each coordinate $z^\mu = x^\mu + iy^\mu$ has a real part ($x^\mu$: the usual spacetime coordinate) and an imaginary part ($y^\mu$: a new ``hidden'' dimension). The real slice $\{Z : \mathrm{Im}(Z) = 0\} \cong \RR^{1,3}$ is ordinary Minkowski spacetime. The imaginary directions encode additional structure (internal degrees of freedom, or an extended spacetime).
\end{remark}

\begin{definition}[Complexified isometry group]\label{def:complex-isom}
The isometry group of $(\CC^{1,3}, \eta_\CC)$ is the \textbf{complex Poincar\'e group}:
\begin{equation}
\Poinc_\CC(c) := \SO(4, \CC) \ltimes \CC^4
\end{equation}
where $\SO(4, \CC) = \{\Lambda \in \GL(4, \CC) : \Lambda^T\eta_\CC\Lambda = \eta_\CC,\; \det\Lambda = 1\}$ is the complexified Lorentz group. The real Poincar\'e group embeds as the real form: $\Poinc(c) \hookrightarrow \Poinc_\CC(c)$.
\end{definition}

\begin{theorem}[Complex kernel]\label{thm:complex-kernel}
$\Ker(\eta_\CC) = \Poinc_\CC(c)$ with $\kappa_\CC = \dim_\CC(\Poinc_\CC) = 10$ (complex dimension) $= 20$ (real dimension). The complexification doubles the number of symmetry generators.
\end{theorem}


% ============================================================================
\section{Twin Hierarchy on Complex Spacetime}
\label{sec:complex-twin-st}
% ============================================================================

\begin{definition}[Componentwise complex twin operations on $\CC^{1,3}$]\label{def:twin-st}
For $Z = (z^0, z^1, z^2, z^3)$ and $W = (w^0, w^1, w^2, w^3)$ in $\CC^{1,3}$:
\begin{equation}
Z \odot_n W := (z^0 \odot_n w^0,\; z^1 \odot_n w^1,\; z^2 \odot_n w^2,\; z^3 \odot_n w^3)
\end{equation}
where each component uses the complex twin operation~\eqref{eq:complex-twin}. This equips $\CC^{1,3}$ with a countably infinite family of group structures $\{(\CC^{1,3}, \odot_n)\}_{n \in \ZZ}$.
\end{definition}

\begin{definition}[Complex twin group]\label{def:twin-group}
The \textbf{complex twin group} $\mathcal{T}_\CC$ is the group generated by all level transitions:
\begin{equation}
\mathcal{T}_\CC := \langle T \rangle, \qquad T: (\CC^{1,3}, \odot_n) \xrightarrow{\sim} (\CC^{1,3}, \odot_{n+1}) \quad \text{for every } n \in \ZZ,
\end{equation}
where $T(Z) = (\exp(z^0), \exp(z^1), \exp(z^2), \exp(z^3))$ is the componentwise complex exponential. By Theorem~\ref{thm:complex-props}(c) the \emph{same} map $T = \exp$ intertwines $\odot_n$ with $\odot_{n+1}$ for all $n$, so $\mathcal{T}_\CC = \langle T \rangle = \{T^{(k)} : k \in \ZZ\}$ is the group of iterates of $\exp$, and $\mathcal{T}_\CC \cong (\ZZ, +)$ via $T^{(k)} \mapsto k$.
\end{definition}

\begin{theorem}[Kernel at each level]\label{thm:kernel-level}
At each level $n$, the kernel (isometry group of $\odot_n$-compatible structures) is:
\begin{equation}
\Ker_n := \{g \in \Aut(\CC^{1,3}) : g(Z \odot_n W) = g(Z) \odot_n g(W)\;\forall Z, W\}
\end{equation}
The kernels at different levels are conjugate: $\Ker_n = T_n \circ \Ker_0 \circ T_n^{-1}$, where $\Ker_0 = \Poinc_\CC(c)$ (the linear isometries at level $0$). Hence all levels have the same kernel dimension ($\kappa_n = 10$ complex).
\end{theorem}


% ============================================================================
\section{Gauge-Gravity Model on $\CC^{1,3}$}
\label{sec:gauge-gravity}
% ============================================================================

\subsection{The total symmetry group}

\begin{construction}[Gauge-gravity group]\label{constr:gauge-gravity}
Define the \textbf{total symmetry group} as the semi-direct product:
\begin{equation}\label{eq:total-group}
\boxed{\mathcal{G} := \underbrace{(\mathcal{T}_\CC \ltimes \Poinc_\CC(c))}_{\text{gravity sector}} \times \underbrace{(\SU(3) \times \SU(2) \times \UU(1))}_{\text{gauge sector}}}
\end{equation}
The gravity sector acts on the \emph{base} $\CC^{1,3}$ (spacetime deformations); the gauge sector acts on the \emph{fiber} (internal symmetries of fields).
\end{construction}

\begin{definition}[Total state space]\label{def:total-state}
The \textbf{total state space} is a fiber bundle:
\begin{equation}
\mathcal{E} = \CC^{1,3} \times_\mathcal{G} F
\end{equation}
where $F$ is the fiber carrying the gauge representation. For the Standard Model: $F = \CC^3 \oplus \CC^2 \oplus \CC$ (color triplet $\oplus$ weak doublet $\oplus$ hypercharge singlet).
\end{definition}

\subsection{Equivariant fields}

\begin{definition}[$(G_{\mathrm{grav}}, G_{\mathrm{gauge}})$-equivariant field]\label{def:field}
A \textbf{physical field} on $\CC^{1,3}$ is a section $\Psi: \CC^{1,3} \to F$ that is equivariant under the total group:
\begin{equation}\label{eq:field-equivariant}
\Psi(g_{\mathrm{grav}} \cdot Z) = \rho(g_{\mathrm{gauge}}(g_{\mathrm{grav}})) \cdot \Psi(Z)
\end{equation}
where $\rho: G_{\mathrm{gauge}} \to \GL(F)$ is the gauge representation and $g_{\mathrm{gauge}}(g_{\mathrm{grav}})$ is the gauge transformation induced by the spacetime transformation via a connection $A$.
\end{definition}

\begin{theorem}[Kernel structure of the total group]\label{thm:total-kernel}
The kernel of the total group $\mathcal{G}$ acting on $\mathcal{E}$ splits as:
\begin{equation}
\Ker(\mathcal{G}) = \Ker_{\mathrm{grav}} \times \Ker_{\mathrm{gauge}}
\end{equation}
where $\Ker_{\mathrm{grav}} = \Poinc_\CC(c)$ (spacetime isometries) and $\Ker_{\mathrm{gauge}} \subseteq \SU(3) \times \SU(2) \times \UU(1)$ (gauge transformations that leave all fields invariant). The deformation space:
\begin{equation}
\Twin(\mathcal{G}) = \underbrace{\mathcal{T}_\CC}_{\text{gravitational deformations}} \times \underbrace{(G_{\mathrm{gauge}}/\Ker_{\mathrm{gauge}})}_{\text{gauge deformations}}
\end{equation}
encodes all physically distinct configurations of gravity + gauge fields.
\end{theorem}

\subsection{The complex Einstein-Yang-Mills system}

\begin{construction}[Field equations from the orbital pattern]\label{constr:eym}
Apply the three-step orbital pattern to the gauge-gravity group~\eqref{eq:total-group}:

\textbf{Step 1 (Symmetry):} The total group $\mathcal{G}$ acts on $\mathcal{E} = \CC^{1,3} \times_\mathcal{G} F$.

\textbf{Step 2 (Conservation):} The equivariance of the field equation under $\mathcal{G}$ implies:
\begin{itemize}
    \item Gravitational conservation: $\nabla_\mu T^{\mu\nu} = 0$ (from $\Poinc_\CC$-equivariance).
    \item Gauge conservation: $D_\mu J^{\mu a} = 0$ (from $G_{\mathrm{gauge}}$-equivariance), where $D_\mu$ is the gauge-covariant derivative and $J^{\mu a}$ is the gauge current.
\end{itemize}

\textbf{Step 3 (Uniqueness):} The unique second-order, $\mathcal{G}$-equivariant, divergence-free equations are:
\begin{align}
G_{\mu\nu}[\eta_\CC^{(n)}] + \Lambda g_{\mu\nu} &= 8\pi T_{\mu\nu} \qquad\text{(complex Einstein, at level $n$)}, \label{eq:complex-einstein} \\
D_\mu F^{\mu\nu a} &= J^{\nu a} \qquad\text{(Yang--Mills)}, \label{eq:ym}
\end{align}
where $\eta_\CC^{(n)}$ is the metric induced by the $\odot_n$ operation and $F^{\mu\nu a}$ is the gauge field strength.
\end{construction}

\begin{remark}[What the complex hierarchy adds]
The complex twin hierarchy adds structure absent from the real case:

\begin{enumerate}[label=(\roman*)]
    \item \textbf{Level $n$ as a new quantum number.} The twin level $n$ parameterizes the ``arithmetic'' of spacetime. Transitions between levels ($n \to n+1$, implemented by the complex exponential) are additional symmetries that have no real analog.
    
    \item \textbf{Branch cuts as topological defects.} The complex logarithm $\log(z)$ has a branch cut along the negative real axis. At level $n \geq 1$, the operation $\odot_n$ involves $\log^{(n)}$, which has $n$ nested branch cuts. These branch cuts are \emph{topological defects} in the complex spacetime --- they divide $\CC^{1,3}$ into sectors where the arithmetic changes discontinuously. They are candidates for domain walls, strings, or other topological objects.
    
    \item \textbf{The imaginary part as internal space.} Writing $z^\mu = x^\mu + iy^\mu$: the real part $x^\mu$ is ordinary spacetime; the imaginary part $y^\mu$ is a 4-dimensional ``internal'' space. The complex twin operations mix $x^\mu$ and $y^\mu$, providing a geometric origin for internal symmetries.
    
    \item \textbf{Analytic continuation and CPT.} The complexified Lorentz group $\SO(4, \CC)$ contains the $PT$ transformation as an element in a different connected component. The CPT theorem (Paper~3b) is the statement that physical amplitudes can be analytically continued through the complexified group, connecting particle and antiparticle representations. The complex spacetime $\CC^{1,3}$ makes this analytic continuation geometric.
\end{enumerate}
\end{remark}


% ============================================================================
\section{The $n$-Level Metric and Deformed Gravity}
\label{sec:n-metric}
% ============================================================================

\begin{definition}[$n$-level metric]\label{def:n-metric}
At level $n$, define the \textbf{$n$-metric} on $\CC^{1,3}$ by transporting $\eta_\CC$ through the gauge map $\phi_n = \log^{(n)}$:
\begin{equation}
\eta_\CC^{(n)}(Z, W) := \eta_\CC(\log^{(n)}(Z), \log^{(n)}(W)) = \sum_\mu \epsilon_\mu \log^{(n)}(z^\mu) \cdot \log^{(n)}(w^\mu)
\end{equation}
where $\epsilon_\mu = (-c^2, +1, +1, +1)$.
\end{definition}

\begin{theorem}[Geodesics at level $n$]\label{thm:n-geodesics}
The geodesics of $\eta_\CC^{(n)}$ are the curves $Z(s)$ satisfying:
\begin{equation}
\frac{d^2}{ds^2}\log^{(n)}(z^\mu(s)) = 0 \qquad \mu = 0, 1, 2, 3
\end{equation}
i.e., $\log^{(n)}(z^\mu(s)) = A^\mu s + B^\mu$ (straight lines in the $\log^{(n)}$-coordinate). In the original coordinate:
\begin{equation}
z^\mu(s) = \exp^{(n)}(A^\mu s + B^\mu)
\end{equation}
At level $0$: straight lines (SR). At level $1$: $z^\mu(s) = e^{A^\mu s + B^\mu}$ (exponential curves). At level $2$: double-exponential curves. Each level has qualitatively different geodesic structure.
\end{theorem}

\begin{theorem}[Force hierarchy from level nesting]\label{thm:force-hierarchy}
The transition from level $n$ to level $n+1$ introduces a ``deformation force'' --- the deviation of $\odot_{n+1}$-geodesics from $\odot_n$-geodesics. The strength of this force scales as:
\begin{equation}
F_{n \to n+1} \sim \frac{d}{ds}\left[\exp^{(n+1)}(A s + B) - \exp^{(n)}(A s + B)\right]
\end{equation}
For large $n$: $\exp^{(n)}$ grows faster than $\exp^{(n-1)}$, so higher levels produce stronger deformation forces. This suggests a hierarchy of forces indexed by the twin level $n$, with gravity ($n = 0 \to 1$) as the weakest and higher levels corresponding to stronger interactions.
\end{theorem}


% ============================================================================
\section{The Full Model: Complex Gauge-Twin-Gravity}
\label{sec:full-model}
% ============================================================================

\begin{construction}[The unified model]\label{constr:full}
The complete model on $\CC^{1,3}$ has:

\textbf{Base manifold:} $\CC^{1,3}$ with the complex twin hierarchy $\{\odot_n\}_{n \in \ZZ}$.

\textbf{Symmetry group:}
\begin{equation}
\mathcal{G}_{\mathrm{full}} = \underbrace{\ZZ}_{\text{level}} \ltimes \underbrace{\Poinc_\CC(c)}_{\text{spacetime}} \times \underbrace{\SU(3)_c \times \SU(2)_L \times \UU(1)_Y}_{\text{gauge}}
\end{equation}

\textbf{Fiber:} $F = \bigoplus_{\text{particles}} V_{(m,s)}^{(\mathbf{r}_c, \mathbf{r}_w, Y)}$ where $V_{(m,s)}$ is the Poincar\'e representation (mass, spin) and $(\mathbf{r}_c, \mathbf{r}_w, Y)$ are the gauge quantum numbers (color, weak isospin, hypercharge).

\textbf{Fields:} Sections $\Psi: \CC^{1,3} \to F$ equivariant under $\mathcal{G}_{\mathrm{full}}$.

\textbf{Equations:}
\begin{enumerate}[label=(\roman*)]
    \item \textbf{Gravity:} Complex Einstein equations~\eqref{eq:complex-einstein} on $(\CC^{1,3}, \eta_\CC^{(n)})$ at level $n$.
    \item \textbf{Gauge:} Yang-Mills equations~\eqref{eq:ym} for the connection on $F$.
    \item \textbf{Matter:} Dirac/Klein-Gordon equations on $\CC^{1,3}$ with gauge coupling.
    \item \textbf{Level coupling:} Transition operators $T_n$ mixing different levels.
\end{enumerate}

\textbf{Real-slice recovery:} Restricting to $\mathrm{Im}(Z) = 0$ and level $n = 0$:
\begin{itemize}
    \item $\CC^{1,3}|_{\mathrm{Re}} = \RR^{1,3}$ (Minkowski).
    \item $\Poinc_\CC|_{\mathrm{Re}} = \Poinc(c)$.
    \item $\eta_\CC^{(0)}|_{\mathrm{Re}} = \eta$ (standard metric).
    \item The model reduces to the Standard Model on flat spacetime.
\end{itemize}
\end{construction}

\begin{theorem}[Consistency of the model]\label{thm:consistency}
The full model satisfies:
\begin{enumerate}[label=(\alph*)]
    \item \textbf{Real-slice consistency:} On $\RR^{1,3} \subset \CC^{1,3}$ at level $0$, all equations reduce to the standard equations of physics (SR, GR, QM, QFT, SM).
    \item \textbf{Level-independence of the kernel:} $\Ker_n \cong \Ker_0 = \Poinc_\CC(c)$ for all $n$ (Theorem~\ref{thm:kernel-level}).
    \item \textbf{Equivariance:} The field equations are $\mathcal{G}_{\mathrm{full}}$-equivariant by construction.
    \item \textbf{Conservation:} The $\Poinc_\CC$-equivariance implies $\nabla_\mu T^{\mu\nu} = 0$; the gauge equivariance implies $D_\mu J^{\mu a} = 0$.
\end{enumerate}
\end{theorem}


% ============================================================================
\section{Discussion and Open Problems}
\label{sec:discussion}
% ============================================================================

\subsection{What is new}

\begin{enumerate}[label=(\roman*)]
    \item The complex twin hierarchy~\eqref{eq:complex-twin} generates a countably infinite family of spacetime arithmetics on $\CC^{1,3}$. Each level $n$ gives a different geodesic structure, a different metric, and a different notion of ``straight line.''
    
    \item The branch cuts of $\log^{(n)}$ introduce topological defects in the complex spacetime that have no analog in the real case. These are candidates for physical objects (domain walls, cosmic strings).
    
    \item The level $n$ provides a new discrete quantum number that indexes the ``arithmetic'' of spacetime. The force hierarchy (Theorem~\ref{thm:force-hierarchy}) suggests that different fundamental forces may correspond to different twin levels.
    
    \item The gauge-gravity crossing~\eqref{eq:total-group} provides a geometric framework where gravity (from the twin hierarchy on the base) and gauge forces (from the fiber) coexist in a single equivariant structure.
\end{enumerate}

\subsection{Open problems}

\begin{enumerate}[label=(\roman*)]
    \item \textbf{Branch cut physics.} What is the physical interpretation of the branch cuts? Do they correspond to observable topological defects? What are the boundary conditions at the branch cuts?
    
    \item \textbf{Level quantization.} Can the twin level $n$ be related to known quantum numbers (baryon number, lepton number, generation)? The discrete nature of $n \in \ZZ$ is suggestive.
    
    \item \textbf{Imaginary dimensions.} What is the physical role of the imaginary parts $y^\mu = \mathrm{Im}(z^\mu)$? Are they compactified (as in Kaluza-Klein)? Are they unobservable (as in complexified GR)?
    
    \item \textbf{Unitarity on $\CC^{1,3}$.} The $\CC$-bilinear metric $\eta_\CC$ is not positive-definite (not even Hermitian). What is the correct notion of ``probability'' on complex spacetime? Is the standard Hilbert space structure adequate, or does a new inner product emerge?
    
    \item \textbf{Predictions.} Does the model predict new particles, new forces, or new phenomena that are absent from the Standard Model + GR? The force hierarchy (Theorem~\ref{thm:force-hierarchy}) suggests higher-level forces, but their strength and coupling remain to be computed.
\end{enumerate}


% ============================================================================
\section{Conclusion}
\label{sec:conclusion}
% ============================================================================

We have presented the theory of twin spaces in three layers:

\textbf{Part~I} (General Theory): The generation mechanism $x \odot_\phi y = \phi^{-1}(\phi(x) \;\alpha\; \phi(y))$ equips any set $(E, \alpha)$ with a family of isomorphic algebraic structures. The kernel-deformation splitting $\mathfrak{g} = \mathfrak{ker} \oplus \mathfrak{def}$ separates symmetries from deformations. Equivariant maps, isomorphism theorems, and orbit-class theorems provide the algebraic infrastructure.

\textbf{Part~II} (Real Spacetime): On $\RR^{1,3}$, the kernel is the Poincar\'e group, the twin operations give Einstein velocity addition, and the deformation space parameterizes gravitational fields (general relativity).

\textbf{Part~III} (Complex Spacetime): On $\CC^{1,3}$, the complex twin hierarchy $z_1 \odot_n z_2 = \exp^{(n)}(\log^{(n)}(z_1) + \log^{(n)}(z_2))$ generates a countably infinite family of spacetime arithmetics. Crossing with the Standard Model gauge group $\SU(3) \times \SU(2) \times \UU(1)$ produces a unified gauge-gravity model where gravity emerges from the twin hierarchy on the base and gauge forces from the fiber. The real slice at level $0$ recovers all known physics.

The complex extension is speculative and exploratory. Its value lies in the questions it raises --- about branch cut physics, level quantization, and the geometry of internal symmetries --- rather than in definitive predictions. It illustrates the power of the twin spaces framework: the same algebraic mechanism that produces Einstein velocity addition on $\RR$ generates an infinite hierarchy of spacetime structures on $\CC$, with the full Standard Model gauge group as a natural companion.


% ============================================================================
\begin{thebibliography}{25}

\bibitem{nembe2026p1} J.~Nemb\'e, ``Special relativity from orbital structure,'' Paper~1 (2026).

\bibitem{nembe2026p2} J.~Nemb\'e, ``General relativity from orbital structure,'' Paper~2 (2026).

\bibitem{nembe2026p3a} J.~Nemb\'e, ``Quantum mechanics from orbital structure,'' Paper~3 (2026).

\bibitem{nembe2026p4} J.~Nemb\'e, ``QFT, the Unifying Observation, and constraints on quantum gravity,'' Paper~4 (2026).

\bibitem{nembe2026vol17} J.~Nemb\'e, ``Twin General Relativity,'' \emph{ROOTS-INSIGHTS} Vol.~17 (2026).

\bibitem{nembe2026vol16} J.~Nemb\'e, ``Twin Information Theory,'' \emph{ROOTS-INSIGHTS} Vol.~16 (2026).

\bibitem{souriau1970} J.-M.~Souriau, \emph{Structure des syst\`emes dynamiques} (Dunod, 1970).

\bibitem{weinberg1995} S.~Weinberg, \emph{The Quantum Theory of Fields}, Vol.~1 (Cambridge Univ.\ Press, 1995).

\bibitem{penrose1967} R.~Penrose, ``Twistor algebra,'' \emph{J.\ Math.\ Phys.}\ \textbf{8}, 345--366 (1967).

\bibitem{newman1973} E.T.~Newman, ``Heaven and its properties,'' \emph{Gen.\ Rel.\ Grav.}\ \textbf{7}, 107--111 (1976).

\bibitem{lebrun1982} C.~LeBrun, ``$\mathcal{H}$-space with a cosmological constant,'' \emph{Proc.\ R.\ Soc.\ A}\ \textbf{380}, 171--185 (1982).

\bibitem{connes1996} A.~Connes, ``Gravity coupled with matter,'' \emph{Comm.\ Math.\ Phys.}\ \textbf{182}, 155--176 (1996).

\bibitem{wald1984} R.M.~Wald, \emph{General Relativity} (Univ.\ Chicago Press, 1984).

\end{thebibliography}

\end{document}
